% =====================================================================
% sec-theory-a.tex — 第 2 节 理论与方法基础（前半部）
% 2.1 概率校准的形式化定义与模型设定（TS 单参数模型与 Platt 双参数对照模型）
% 2.2 极大似然估计、Score 方程解析推导、Fisher 信息量与严格凹性数学证明
% =====================================================================
\section{理论与方法基础}\label{sec:theory}

\subsection{概率校准的形式化定义与模型设定}

在多分类机器学习与统计模式识别框架下，设样本观测值由多元随机向量 $X \in \mathcal{X} \subset \mathbb{R}^d$ 与真实离散类别标签 $Y \in \mathcal{Y} = \{1, 2, \dots, C\}$ 构成，其联合概率测度服从未知总体分布 $P(X, Y)$。深度分类模型的前向推理网络可形式化为一个可微参数映射 $f: \mathcal{X} \to \mathbb{R}^C$。对于输入特征向量 $X_i$，网络在最终线性分类层输出未经归一化的实值对数几率（Logit）向量：
\begin{equation}\label{eq:logit_vector}
\boldsymbol{z}_i = f(X_i) = \big(z_{i1}, z_{i2}, \dots, z_{iC}\big)^\top \in \mathbb{R}^C.
\end{equation}

在经典标准架构中，模型输出概率采用标准 Softmax 激活函数进行映射，得到经验后验预测分布向量 $\hat{\boldsymbol{p}}_i = (\hat p_{i1}, \dots, \hat p_{iC})^\top$：
\begin{equation}\label{eq:standard_softmax}
\hat p_{ic} = \frac{\exp(z_{ic})}{\sum_{k=1}^C \exp(z_{ik})}, \quad c = 1, 2, \dots, C.
\end{equation}
离散分类判决依据贝叶斯最大后验概率（MAP）准则给出：
\begin{equation}\label{eq:map_decision}
\hat Y_i = \operatorname{argmax}_{c \in \{1,\dots,C\}} \hat p_{ic} = \operatorname{argmax}_{c \in \{1,\dots,C\}} z_{ic},
\end{equation}
对应的最大预测置信度记为 $\hat P_i = \max_{c} \hat p_{ic} = \hat p_{i, \hat Y_i} \in [1/C, 1]$。

\begin{defn}[条件概率校准性]\label{defn:calibration}
若对于任意给定的概率预测水平 $p \in [0, 1]$，真实标签与模型判决类别一致的条件概率精确恒等于该预测概率水平，即：
\begin{equation}\label{eq:calibration_formal}
\mathbb{P}\big(Y = \hat Y \;\big|\; \hat P = p\big) = p, \quad \forall p \in [0, 1],
\end{equation}
则称该多分类模型满足严格条件概率校准。
\end{defn}

然而，在由大量卷积层与残差连接构成的现代网络中，经验损失最小化训练导致对数几率向量 $\boldsymbol{z}_i$ 的模长显著过大，打破了式~\eqref{eq:calibration_formal}的平衡，导致系统性的过度自信。为消除这一概率失配，我们建立后处理参数校准统计模型。

\subsubsection{单参数 Temperature Scaling 模型设定}

Temperature Scaling（温度缩放，TS）是一种作用于网络最终对数几率向量的多分类单参数后处理变换。本文将其形式化为单参数指数族分布。引入非负连续温度标量参数 $T \in (0, \infty)$，或等价引入其代数倒数——\textbf{逆温度参数}（Inverse Temperature Parameter）$\beta \equiv 1/T \in (0, \infty)$。给定样本 $i$ 的对数几率向量 $\boldsymbol{z}_i$，校准后的模型条件后验概率分布定义为：
\begin{equation}\label{eq:ts_model}
\begin{split}
p_{ic}(\beta) \equiv \mathbb{P}(Y_i = c \mid \boldsymbol{z}_i; \beta) &= \frac{\exp(\beta z_{ic})}{\sum_{k=1}^C \exp(\beta z_{ik})} \\
&= \operatorname{Softmax}_c(\beta \boldsymbol{z}_i), \quad c = 1, 2, \dots, C.
\end{split}
\end{equation}

单参数 Temperature Scaling 具有两个至关重要的数理统计特征：
\begin{prop}[严格保序性（Monotonic Order Preservation）]\label{prop:order_preservation}
对于任意给定的有限常数 $\beta \in (0, \infty)$ 以及任意样本对数几率分量，若存在两类别分量 $c, k \in \{1, \dots, C\}$ 满足 $z_{ic} > z_{ik}$，则当且仅当 $\beta > 0$ 时恒有：
\begin{equation}\label{eq:order_inequality}
p_{ic}(\beta) > p_{ik}(\beta).
\end{equation}
特别地，校准后的最大后验分类决策与未校准的原始网络决策严格恒等：
\begin{equation}\label{eq:argmax_invariance}
\operatorname{argmax}_{c \in \{1,\dots,C\}} p_{ic}(\beta) \equiv \operatorname{argmax}_{c \in \{1,\dots,C\}} z_{ic} \equiv \hat Y_i.
\end{equation}
\end{prop}

\begin{proof}
由指数函数的严格单调递增性，对任意 $\beta > 0$，有 $z_{ic} > z_{ik} \iff \beta z_{ic} > \beta z_{ik} \iff \exp(\beta z_{ic}) > \exp(\beta z_{ik})$。由于式~\eqref{eq:ts_model}中各分母归一化项 $\sum_k \exp(\beta z_{ik})$ 恒为严格正实数且与类别指标 $c$ 无关，两端同除分母不等号方向保持不变，即证式~\eqref{eq:order_inequality}。进而取极大值算子即可直接得出式~\eqref{eq:argmax_invariance}恒等成立。
\end{proof}

命题~\ref{prop:order_preservation}在理论上确立了 Temperature Scaling 的保序性质：\textbf{在后处理优化概率校准度的全过程中，模型的离散 Top-1 预测准确率在理论上保持不变，不会产生判别性能退化}。
从参数几何意义来看，当 $\beta = 1$（即 $T = 1$）时，模型精确退化为原始未校准网络；当 $\beta \in (0, 1)$（即 $T > 1$）时，$\beta$ 线性压缩了各类别 logit 之间的欧氏间距，使 Softmax 概率分布向中心平滑回缩，从而系统性抑制高置信度；当 $\beta \to 0$（即 $T \to \infty$）时，后验分布趋于离散均匀分布 $(1/C, \dots, 1/C)^\top$。

\subsubsection{双参数 Platt Scaling 对照模型设定}

为深入评估单参数缩放的简约性与相对优势，本文引入经典的双参数 Platt Scaling（Platt, 1999）作为基准对照模型。Platt Scaling 建立在模型 Top-1 预测的对数几率（Log-odds）标量变换之上。定义样本 $i$ 的标量预测几率为：
\begin{equation}\label{eq:log_odds}
x_i \equiv \operatorname{logit}(\hat P_i) = \log\left(\frac{\hat P_i}{1 - \hat P_i}\right) \in \mathbb{R}.
\end{equation}
在二项逻辑斯蒂回归（Logistic Regression）框架下，校准模型设定两个自由实数参数——偏置截距项 $\alpha \in \mathbb{R}$ 与斜率缩放项 $\beta_{\mathrm{platt}} \in \mathbb{R}$，重新估计 Top-1 类别预测为真实的概率 $q_i$：
\begin{equation}\label{eq:platt_formula}
\begin{split}
q_i(\alpha, \beta_{\mathrm{platt}}) &\equiv \mathbb{P}\big(Y_i = \hat Y_i \;\big|\; x_i; \alpha, \beta_{\mathrm{platt}}\big) \\
&= \sigma(\alpha + \beta_{\mathrm{platt}} x_i) = \frac{1}{1 + \exp\big(-(\alpha + \beta_{\mathrm{platt}} x_i)\big)},
\end{split}
\end{equation}
其中 $\sigma(\cdot)$ 为标准单变量 Sigmoid 激活函数。对于其余非 Top-1 类别的剩余概率，采用条件独立比例重归一化：
\begin{equation}\label{eq:platt_renorm}
\tilde p_{ic} = \frac{1 - q_i}{1 - \hat P_i} \hat p_{ic}, \quad \forall c \ne \hat Y_i.
\end{equation}

与单参数 Temperature Scaling 相比，Platt Scaling 引入了无约束的截距参数 $\alpha$。当 $\alpha \ne 0$ 时，校准映射不再通过原点，导致预测分类边界发生全局刚性平移。如果某样本的校准后概率满足 $q_i < 1/C$ 或引起非 Top-1 类别概率倒挂，则可能导致样本分类标签翻转，从而对模型的原始 Top-1 准确率造成外生性扰动。

\subsection{极大似然估计、Score 方程与严格凹性证明}

设有一组规模为 $n$ 的独立同分布（i.i.d.）校准数据集 $\mathcal{D}_{\mathrm{cal}} = \{(\boldsymbol{z}_i, y_i)\}_{i=1}^n$，其中 $\boldsymbol{z}_i \in \mathbb{R}^C$ 为预训练网络输出的对数几率向量，$y_i \in \{1, \dots, C\}$ 为真实的离散类别标签。引入独热（One-Hot）指示变量向量 $\boldsymbol{y}_i = (y_{i1}, \dots, y_{iC})^\top \in \{0, 1\}^C$，其中 $y_{ic} = \mathbf{1}(y_i = c)$。

\subsubsection{对数似然函数与解析 Score 方程}

在多分类指数族似然准则下，样本条件联合对数似然函数（Log-Likelihood Function）表示为逆温度参数 $\beta$ 的一元显式函数：
\begin{align}
\ell(\beta) &\equiv \sum_{i=1}^n \sum_{c=1}^C y_{ic} \log p_{ic}(\beta) \nonumber\\
&= \sum_{i=1}^n \sum_{c=1}^C y_{ic} \left[ \beta z_{ic} - \log \left( \sum_{k=1}^C \exp(\beta z_{ik}) \right) \right] \nonumber\\
&= \sum_{i=1}^n \left[ \beta z_{i, y_i} - \log \left( \sum_{k=1}^C \exp(\beta z_{ik}) \right) \right]. \label{eq:full_loglik}
\end{align}

对对数似然函数 $\ell(\beta)$ 关于参数 $\beta$ 求取一阶解析导数，可得统计学中的 \textbf{Score 函数}（记为 $U(\beta)$）：
\begin{align}
U(\beta) &\equiv \frac{\mathrm{d}\ell(\beta)}{\mathrm{d}\beta} = \sum_{i=1}^n \left[ z_{i, y_i} - \frac{\sum_{k=1}^C z_{ik} \exp(\beta z_{ik})}{\sum_{m=1}^C \exp(\beta z_{im})} \right] \nonumber\\
&= \sum_{i=1}^n \left[ z_{i, y_i} - \sum_{c=1}^C p_{ic}(\beta) z_{ic} \right]. \label{eq:score_derivation}
\end{align}

定义离散概率测度 $p_i(\beta)$ 下，样本 $i$ 处对数几率随机变量 $Z_i \in \{z_{i1}, \dots, z_{iC}\}$ 的条件期望为：
\begin{equation}\label{eq:conditional_expectation}
\mathbb{E}_{p_i(\beta)}[Z_i] \equiv \sum_{c=1}^C p_{ic}(\beta) z_{ic}.
\end{equation}
则 Score 函数式~\eqref{eq:score_derivation}可改写为真实观测值与模型期望值的残差和：
\begin{equation}\label{eq:score_clean}
U(\beta) = \sum_{i=1}^n \Big( z_{i, y_i} - \mathbb{E}_{p_i(\beta)}[Z_i] \Big).
\end{equation}
极大似然估计方程（Likelihood Equation）要求置 Score 函数为零，即 $U(\hat\beta) = 0$。该解析式表明：\textbf{在 MLE 估计点 $\hat\beta$ 处，全体校准样本中真实所属类别的 logit 总和等于在校准后模型分布测度下所计算的期望 logit 总和}。

\subsubsection{二阶导数、Fisher 信息量与严格凹性证明}

进一步对 Score 函数 $U(\beta)$ 求取关于参数 $\beta$ 的二阶导数。注意到 Softmax 概率关于 $\beta$ 的偏导数为：
\begin{equation}\label{eq:softmax_derivative}
\begin{split}
\frac{\mathrm{d} p_{ic}(\beta)}{\mathrm{d}\beta} &= p_{ic}(\beta) z_{ic} - p_{ic}(\beta) \left( \sum_{k=1}^C p_{ik}(\beta) z_{ik} \right) \\
&= p_{ic}(\beta) \big( z_{ic} - \mathbb{E}_{p_i(\beta)}[Z_i] \big).
\end{split}
\end{equation}
利用式~\eqref{eq:softmax_derivative}，直接对式~\eqref{eq:score_derivation}求导：
\begin{align}
\frac{\mathrm{d}^2\ell(\beta)}{\mathrm{d}\beta^2} &= \frac{\mathrm{d} U(\beta)}{\mathrm{d}\beta} = - \sum_{i=1}^n \sum_{c=1}^C z_{ic} \frac{\mathrm{d} p_{ic}(\beta)}{\mathrm{d}\beta} \nonumber\\
&= - \sum_{i=1}^n \sum_{c=1}^C z_{ic} \left[ p_{ic}(\beta) z_{ic} - p_{ic}(\beta) \sum_{k=1}^C p_{ik}(\beta) z_{ik} \right] \nonumber\\
&= - \sum_{i=1}^n \left[ \sum_{c=1}^C p_{ic}(\beta) z_{ic}^2 - \left( \sum_{c=1}^C p_{ic}(\beta) z_{ic} \right)^2 \right]. \label{eq:second_deriv}
\end{align}

根据数理统计理论，定义样本\textbf{观测 Fisher 信息量}（Observed Fisher Information）为负二阶导数：
\begin{equation}\label{eq:fisher_defn}
I_n(\beta) \equiv - \frac{\mathrm{d}^2\ell(\beta)}{\mathrm{d}\beta^2} = - \ell''(\beta).
\end{equation}
式~\eqref{eq:second_deriv}中方括号内的第二项为条件期望的平方，第一项为条件二阶矩 $\mathbb{E}_{p_i(\beta)}[Z_i^2] = \sum_{c=1}^C p_{ic}(\beta) z_{ic}^2$。因此，该方括号恒等于离散随机变量 $Z_i$ 在分布 $p_i(\beta)$ 下的\textbf{条件方差}（Conditional Variance）：
\begin{equation}\label{eq:fisher_variance}
I_n(\beta) = \sum_{i=1}^n \operatorname{Var}_{p_i(\beta)}(Z_i).
\end{equation}

基于式~\eqref{eq:fisher_variance}，我们给出本模型最核心的数学收敛性定理：

\begin{thm}[对数似然的严格凹性与 MLE 唯一全局最优性]\label{thm:strict_concavity}
设校准集 $\mathcal{D}_{\mathrm{cal}}$ 中至少存在一个样本 $i$，其输出对数几率向量是非退化的（即存在至少两类 $c \ne k$ 使得 $z_{ic} \ne z_{ik}$）。则对于任意参数值 $\beta \in (0, \infty)$，观测 Fisher 信息量严格正定：
\begin{equation}\label{eq:fisher_positive}
I_n(\beta) > 0, \quad \forall \beta \in (0, \infty).
\end{equation}
进而，对数似然函数 $\ell(\beta)$ 在开区间 $(0, \infty)$ 上处处严格凹（Strictly Concave）。若 Score 方程 $U(\beta) = 0$ 在定义域内存在实根，则极大似然估计解 $\hat\beta$ 全局唯一。
\end{thm}

\begin{proof}
根据概率论基本性质，对于任意非退化离散随机变量，其方差恒满足 $\operatorname{Var}(X) \ge 0$，且取等号当且仅当随机变量以概率 1 退化为常数。
对于任意样本 $i$ 以及任意有限参数 $\beta \in (0, \infty)$，由指数映射性质可知 Softmax 条件概率恒满足严格内点条件：$p_{ic}(\beta) \in (0, 1)$。由于假定存在样本 $i$ 使得类别间 logit 取值不完全相等，该样本对数几率在分布测度 $p_i(\beta)$ 下不退化为常数，因此该样本的条件方差严格大于零：
\[
\operatorname{Var}_{p_i(\beta)}(Z_i) > 0.
\]
将其代入式~\eqref{eq:fisher_variance}，作为非负项与至少一个严格正项的有限求和，必然得出：
\[
I_n(\beta) = \sum_{i=1}^n \operatorname{Var}_{p_i(\beta)}(Z_i) > 0, \quad \forall \beta \in (0, \infty).
\]
从而对数似然的二阶导数处处严格为负：
\[
\ell''(\beta) = -I_n(\beta) < 0, \quad \forall \beta \in (0, \infty).
\]
由凸分析基本定理，定义在凸集上的二次连续可微函数若其二阶导数恒严格负，则该函数在该凸集上严格凹。若严格凹函数在开区间内部存在一阶导数为零的驻点 $\hat\beta$（即 $U(\hat\beta)=0$），则该点必为唯一的全局极大值点。定理证毕。
\end{proof}

\subsubsection{阻尼 Newton--Raphson 求解}

得益于定理~\ref{thm:strict_concavity}所保证的全局严格凹性，对数似然曲面上不存在局部非凸陷阱或伪驻点。此外，二阶导数式~\eqref{eq:fisher_variance}具备条件方差的闭式表达，便于采用二阶牛顿法高效求解。本文采用标准的带回溯线搜索（Backtracking Line Search）的阻尼 Newton--Raphson 迭代法求解 $\hat\beta$：
\begin{equation}\label{eq:newton_iteration}
\beta^{(t+1)} = \beta^{(t)} + \gamma_t \frac{U(\beta^{(t)})}{I_n(\beta^{(t)})}, \quad t = 0, 1, 2, \dots
\end{equation}
其中，初始点设为未校准状态 $\beta^{(0)} = 1.0$；$\gamma_t \in (0, 1]$ 为阻尼步长因子，初始设为 $1.0$。在每次迭代中，若试探步无法满足单调似然上升条件：
\begin{equation}\label{eq:armijo}
\ell\left(\beta^{(t)} + \gamma_t \frac{U(\beta^{(t)})}{I_n(\beta^{(t)})}\right) < \ell\left(\beta^{(t)}\right),
\end{equation}
则按比例收缩步长 $\gamma_t \leftarrow \gamma_t / 2$；同时对更新值施加保护截断 $\beta \ge 0.05$ 以防止边界发散。迭代收敛判据设定为连续两次更新步长的绝对变化量小于极小公差：$|\beta^{(t+1)} - \beta^{(t)}| < 10^{-12}$。在实测中，得益于二阶导数的精确解析性，该算法通常在 $8 \sim 15$ 次迭代内即可达成机器精度收敛。
